\section{Discussion}

The results reframe an earlier, apparently negative finding. A first pass of this work had reported that
neither representation could rank cell lines, with out-of-fold correlations near zero over the full drug
panel. That outcome was substantially an artifact of the training target rather than a statement about the
data or the representation. The initial score, a dose-averaged viability, left the per-drug variances
untouched, so in the shared multi-task loss a handful of high-variance and largely uninformative
compounds dominated the representation. Standardizing the target per drug removes this imbalance, and with
it the model ranks unseen cell lines at correlations of \rhoFullScgpt{} across the panel and up to
\rhoTenScgpt{} on a learnable subset. The standardization is best understood as a necessary preprocessing
choice for a heterogeneous multi-task objective, not as a scientific result in itself; a head-to-head
comparison of the candidate scores is informative mainly because it shows that the curve fit alone changes
nothing and the standardization changes everything. The dominance of the target change is direct: on the
unstandardized targets the same model collapses at full head count (Table~\ref{tab:target}), and no other
single intervention on the broken configuration approaches it (Figure~\ref{fig:rescue}).

The corrected picture is modest but genuine, and it is best read against the standards the DrEval study
sets for the field. The model clears the NaiveMeanEffects baseline that about half of published models
fail to beat, and it edges DrEval's own reference models on identical features, but it remains below the best
model that study reports and shows only a faint representation advantage once the mean effects are
normalized away. The most informative internal result points the same way: a ridge regression on
cell-line mean embeddings matches the deep model, so the information the current pipeline exploits is
largely captured by a per-line average. This is the cell-line-effect bias that DrEval was designed to
expose, reproduced here independently. The single-cell resolution and the scGPT embedding together add a
small margin over that average, and scGPT needs a nonlinear head to realize it, which is the first
concrete evidence that the embedding carries structure a linear model cannot reach. This lower reliance on
memorization is the effect the guiding hypothesis anticipated, and it persists on the corrected target,
where the scGPT model's training-to-validation gap (\gapScgpt{}) is far smaller than that of PCA
(\gapPca{}) under an identical trunk.

Taken together, the evidence indicates that the limiting factor at this stage is the label side rather
than the model side. The task is defined on roughly \NIndep{} independent cell lines, each carrying one
bulk value broadcast onto hundreds of cells, and model-side tuning is exhausted across regularization,
capacity and batch size. Improving performance therefore requires acting on the labels and the data, not
on the architecture.
